\subsection{实射影空间的拓扑}

我们回忆代数中的下述定义：给定一个除环或域 \(\mathbf{K}\) ，把过 \(o\) 的直线（即一维向量子空间）——左向量空间 \(\mathbf{K}_s^{n+1}\) （\(\mathbf{K}\) 上的）中的——所成的集，称为 \(n\) 维左射影空间（\(\mathbf{K}\) 上的），记作 \(\mathbf{P}_n(\mathbf{K})\)。

若把 \(K_{s}^{n+1}\) 中过 o 的每条直线对应于同一条去掉原点的直线，我们就得到 \(\mathbf{P}_{n}(\mathbf{K})\) 到 \(K_{n+1}^{*}\) （\(\{o\}\) 在 \(K_{s}^{n+1}\) 中的补集）关于下述等价关系 \(\Delta_{n}(\mathbf{K})\) 的商上的一个双射，这个等价关系定义在 \(K_{n+1}^{*}\) 的向量 x 与 y 之间：“存在 \(t \in K\) 使得 \(t \neq o\) 且 y = tx”。在下文中我们将把 \(P_{n}(\mathbf{K})\) 与这个商集等同起来。在射影空间的理论中，我们取 N 的区间 \([o, n]\) 作为 \(K_{n+1}^{*}\) 的点的坐标的指标集。坐标 \(x_{i}\) （ \(o \leqslant i \leqslant n\) ）——\(K_{n+1}^{*}\) 中典范象为 \(x \in P_{n}(\mathbf{K})\) 的任何一点的坐标——组成所谓的点 x 的齐次坐标系。

对每个整数 \(p\) （使 \(-1 \leqslant p \leqslant n\) 的），\(\mathbf{P}_n(\mathbf{K})\) 中一个 \(p + 1\) 维向量子空间（\(\mathrm{K}_s^{n+1}\) 的，去掉原点）的典范象称为 \(p\) 维射影线性簇。\(p\) 个点（\(\mathbf{P}_n(\mathbf{K})\) 中的）组成的一个系称为自由的，如果它由 \(p\) 个点——\(\mathrm{K}_{n+1}^*\) 的、在向量空间 \(\mathrm{K}_s^{n+1}\) 中构成自由系的——的典范象组成。\(\mathbf{P}_n(\mathbf{K})\) 中由 \(p + 1\) 个点的自由系生成的射影线性簇（即包含这 \(p + 1\) 个点的最小射影线性簇）是 \(p\) 维的。

当 K 是域 R 时，相应的射影空间可以赋以拓扑，我们将研究的就是这些空间。

定义 1. 射影空间 \(\mathbf{P}_n(\mathbf{R})\) ——赋以 \(\mathbf{R}_{n+1}^*\) 的拓扑关于等价关系 \(\Delta_n(\mathbf{R})\) 的商拓扑——，称为 \(n\) 维实射影空间。

射影空间 \(\mathbf{P}_{1}(\mathbf{R})\) 称为实射影直线，\(\mathbf{P}_{2}(\mathbf{R})\) 称为实射影平面。

每当没有混淆的危险时，我们将用 \(P_{n}\) 与 \(\Delta_{n}\) 代替 \(\mathbf{P}_{n}(\mathbf{R})\) 与 \(\Delta_{n}(\mathbf{R})\)。

命题 1. 射影空间 \(P_{n}\) 是豪斯多夫的。

我们首先证明关系 \(\Delta_{n}\) 是开的（第 I 章，§ 5，第 2 小节）。设 A 是 \(\mathbf{R}_{n+1}^{*}\) 中的开集；要使 A 关于 \(\Delta_{n}\) 饱和，须取与 A 位似的诸集 \(t\mathbf{A}\) 的并，这里 \(t\) 取遍 \(\neq 0\) 的实数的集；由于这些集中每一个都是开的，它们的并也是开的。

由第 I 章，§ 8，第 3 小节的命题 8，只要证明 \(\mathbf{R}_{n+1}^{*}\times\mathbf{R}_{n+1}^{*}\) 的由关系 \(\Delta_{n}\) 定义的子集 M 是闭的，命题即得证明。设 \((x,y)\) 是 \(\mathbf{R}_{n+1}^{*}\times\mathbf{R}_{n+1}^{*}\) 中属于 M 的闭包的一点。若 \(x=(x_{i})\) ，则有指标 \(i\) 使 \(x_{i}\neq0\) ；因此存在 \((x,y)\) 的邻域 V，使得对每一点 \((x^{\prime},y^{\prime})\in\mathbf{M}\cap\mathbf{V}\) ，第 \(i\) 个坐标 \(x_{i}^{\prime}\) （\(x^{\prime}\) 的）不是 o。当 \((x^{\prime},y^{\prime})\) 保持属于 M 而趋于 \((x,y)\) 时，\(y_{i}^{\prime}x_{i}^{\prime-1}\) 趋于 \(t=y_{i}x_{i}^{-1}\) ；由于 \(y^{\prime}=(y_{i}^{\prime}x_{i}^{\prime-1})x^{\prime}\) ，取极限便见 \(y=tx\) ，这表明 \((x,y)\in\mathbf{M}\) 。

命题 2. 射影空间 \(\mathbf{P}_{n}\) 是紧的且连通的，并且同胚于球面 \(\mathbf{S}_{n}\) 关于 \(\Delta_{n}\) 在该球面上诱导的等价关系的商。

令 \(\Delta_n'\) 是 \(\mathbf{S}_n\) 上由 \(\Delta_n\) 诱导的等价关系（\(\Delta_n'\) 的等价类是 \(\mathbf{S}_n\) 的对径点对）。映射 \(x \to x/||x||\) （\(\mathbf{R}_{n+1}^*\) 到 \(\mathbf{S}_n\) 上的）是连续的，因此（引理 2）\(\mathbf{P}_n\) 同胚于 \(\mathbf{S}_n/\Delta_n'\)。由于 \(\mathbf{S}_n\) 是紧的且连通的，\(\mathbf{S}_n\) 的每个豪斯多夫商空间也是紧的且连通的（第 I 章，§ 9，第 4 小节，定理 2 的推论 1；§ 11，第 3 小节，命题 7）。

命题 3. 若 \(n \geqslant 0\) ，射影空间 \(\mathbf{P}_n\) 同胚于把球 \(\mathbf{B}_n\) 中 \(\mathbf{S}_{n-1}\) 的每个点与其对径点等同起来所得到的商空间。

设 H 是 \(S_{n}\) 的由 \(x_{0} \leqslant 0\) 定义的闭半球面。\(P_{n}\) ——它同胚于 \(S_{n}\) 关于关系 \(\Delta_{n}^{\prime}\) 的商空间——也同胚于 \(S_{n}\) 的子集 H 关于关系 \(\Delta_{n}^{\prime\prime}\) （由 \(\Delta_{n}^{\prime}\) 在 H 上诱导的）的商。由于 \(\Delta_{n}^{\prime}\) 的每个等价类都至少与 H 交于一点，因此（引理 1）只需验证：关于 \(\Delta_{n}^{\prime}\) 把 U （H 的开子集，关于 \(\Delta_{n}^{\prime\prime}\) 饱和）加以饱和化，就得到 V ——\(S_{n}\) 的一个开子集。现在，若 \(a = (a_{i}) \in U\) 且 \(a_{0} < 0\) ，则存在邻域 W （a 的，在 \(S_{n}\) 中的）含于 U，而 W 与 \(-W\) 的并是 a 的一个关于 \(\Delta_{n}^{\prime}\) 饱和且含于 V 的邻域。反之若 \(a_{0} = 0\) ，我们有 \(- a \in U\) ，且存在 r \textgreater{} o，使得点 \(x \in H\) ——满足关系 \(||x-a||<r, ||x+a||<r\) 之一的——所成的集合含于 U；点 \(x\in S_{n}\) ——满足这两个关系之一的——所成的集合就是 a 的那个关于 \(\Delta_{n}^{\prime}\) 饱和且含于 V 的邻域。

注意，商空间 \(H/\Delta_{n}^{\prime\prime}\) 是在 H 中把 \(S_{n-1}\) （即 H 与超平面 \(x_{0}=0\) 的交）的每个点与其相反点等同起来而得到的。要完成证明，只需注意以 \(e_{0}\) 为顶点的球极投影（§ 2，第 4 小节）是 H 到 \(B_{n}\) 上的同胚，且保持 \(S_{n-1}\) 的点不变。

